\section{Complexity}
\label{sec:complexity}

This section bounds the cost of deciding entailment. The input lists the
arity $n$, the variance table $c(0),\ldots,c(n-1)$, the variable set $V$, the
query, and the literals of $\phi$, each constructor literal listing all $n$
child variables. Write $\insize\phi$ for its size; thus $n\le\insize\phi$ and
$k=|V|\le\insize\phi$. The readers of \cref{sec:label-reader} have
$N=2\cdot4^{2k}=2^{4k+1}$ states. The constants hidden in $O$-notation below
are absolute: they depend neither on $n$, nor on $c$, nor on the domain. We
write 2-EXPSPACE for the problems decidable in deterministic space
$2^{2^{p(\insize\phi)}}$, and 3-EXPTIME for those decidable in deterministic
time $2^{2^{2^{p(\insize\phi)}}}$, for polynomials $p$. Deterministic
classes, such as these and PSPACE, are closed under complement, so their
bounds hold for entailment and non-entailment alike.

This section concentrates on constructors of arity $n\ge2$.
\Cref{sec:complexity-upper} proves an upper bound that holds for every
arity: a search that stores one level of the block hierarchy at a time places
entailment in 2-EXPSPACE and in 3-EXPTIME. For $n\le1$ stronger bounds hold,
PSPACE at arity one and polynomial time at arity zero
(\cref{app:small-arity}). \Cref{sec:complexity-lower} proves lower bounds for
a covariant constructor of each fixed arity $n\ge2$: shortest holes can be
exponentially long in $k$, and entailment is PSPACE-hard. All these bounds
hold in each of the three domains.

\subsection{Upper Bounds}
\label{sec:complexity-upper}

\paragraph{Hole Length}
The constants of \eqref{eq:bounds} are $B=4N-2$, $d=8N^2-2N-1$, and
$H=\holeBound(N,d)$. By \cref{thm:hole-decision,lem:elementary-estimate},
\[
 H\le2^{2^{1280(N+1)^5}},
\]
a triply exponential bound $2^{2^{2^{O(k)}}}$ in $k$. When $\phi$ is
satisfiable in the domain considered and the entailment fails, a hole of
length at most $H$ of one of the two label readers certifies the failure
(\cref{thm:holes,prop:transfer}). It is an unsafe word, and it can be checked
in time polynomial in its length and in $\insize\phi$. The bound concerns the
length of an unsafe word, not the node count of a fully expanded
countermodel. Enumerating all words up to length $H$, as \cref{fig:algorithm}
does, examines $\sum_{L\le H}n^L$ words per reader, four exponentials in $k$
for fixed $n\ge2$. The summary search below decides the same question without
storing a word.

\paragraph{Polynomial-Time Tests}
Every test of the reduction other than the search for holes is cheap.

\begin{lemma}[Polynomial-time tests]\label{lem:poly-tests}
From a flat system $\phi$ with $k$ variables and the query, the following are
computable in time polynomial in $\insize\phi$: the signed system
$\psi=\phi^\pm$ and its closure; whether it has a
label clash or a cycle clash; and, for either label reader and states given
as bit vectors, the initial state, the successor of a state under a letter,
whether a state is rejected, and whether a pair of states is admitted. In
particular, $\Sat_{\mathcal D,c}(\phi)$ is decidable in polynomial time for
each of the three domains $\mathcal D$.
\end{lemma}

The closure has at most $(2k)^2$ pairs and is computed by saturation. Every
reader operation is a set operation on the $2k$ signed variables, computed
from the closure. The clash tests decide satisfiability by \eqref{eq:g19},
\cref{thm:sat}, \eqref{eq:sat-reg}, and \cref{lem:finite-signed,thm:finite-sat}.
The proof is in \cref{app:complexity}.

\paragraph{Searching Without Storing the Word}
\Cref{thm:hole-decision} bounds the length of a shortest hole, but a hole of
that length takes triply exponential space to write down. The search below
stores much less. It rests on three observations about \cref{sec:search}.
First, a lettered monitor with $N$ states and a nonempty hole of depth at
most $d$ has such a hole whose hierarchy uses at most $\DictBound(N,d)$ distinct letters at each positive
level (\cref{lem:dictionary-bound}). The $h$ blocks of a level, over a monitor
with $N+j$ states, can be shortened together below $\mathit{col}(N+j,h)$
letters without changing their summary, and a unary hole of such a monitor
can be shortened to at most $N+j$ letters. Second, compression cuts a given
hole top-down, but the derived monitors are built bottom-up:
$\mathcal M_{j+1}$ is the derivative of $\mathcal M_j$ at the marker
$\ell_j$. Restricted to the blocks that are the letters of level $j+1$, and
up to their names, it depends only on the summary of these blocks and on the
targets of $\mathcal M_j$ at the marker, not on the blocks themselves
(\cref{lem:summary-monitor}). A search can therefore forget the blocks and
call them $1,\ldots,h$. Third, the summary is computed by an automaton that
reads the blocks column by column (\cref{fig:columns}), so the columns can be
guessed and discarded one at a time.

The search thus keeps only the tables of a current monitor over the abstract
letters $1,\ldots,h$. It guesses the next level's blocks column by column,
keeps their running summary, and accepts when a short unary word is a hole of
the current monitor (\cref{fig:summary-search,fig:summary-levels}). The stored
tables are doubly exponential in $k$, although the hole they describe may be
triply exponentially long.

\begin{figure}[tbp]
\small
\begin{minipage}{0.94\linewidth}
\textbf{Input:} a lettered monitor $\mathcal M_0$ with $N$ states over a
finite alphabet $\Gamma$, given by its tables, and a depth $d$.\\
\textbf{Stored:} a level $j\le d$ and a current monitor $\mathcal W$ with
$N+j$ states over a current alphabet $\Gamma_{\mathcal W}$; initially $j=0$,
$\mathcal W=\mathcal M_0$, and $\Gamma_{\mathcal W}=\Gamma$.\\
\textbf{Repeat} one of two nondeterministic moves:
\begin{enumerate}[leftmargin=*,itemsep=2pt]
\item \emph{Accept.} Guess a letter $z\in\Gamma_{\mathcal W}$ and a length
$1\le m'\le N+j$. Accept if $z^{m'}$ is a hole of $\mathcal W$.
\item \emph{Derive}, if $j<d$. Guess a marker $\ell\in\Gamma_{\mathcal W}$ and
a number $1\le h\le\DictBound(N,d)$. Guess a valid column word with $h$
tracks over $\Gamma_{\mathcal W}\setminus\{\ell\}$ and fewer than
$\mathit{col}(N+j,h)$ columns, one column at a time. After each column,
update the summary $\Theta$ of the tracks with respect to $\mathcal W$ and
$\ell$, and discard the column. Stop without accepting unless $\Theta$
declares the $h$ blocks pairwise distinct. Then replace $\mathcal W$ by the
monitor $\mathcal W[\ell,\Theta]$ over $\Gamma_{\mathcal W}=[h]$ that the
summary determines, and $j$ by $j+1$.
\end{enumerate}
\end{minipage}
\caption{The summary search of \cref{prop:summary-search}. Here
$[h]=\{1,\ldots,h\}$, and $\mathcal W[\ell,\Theta]$ has one new state
(\cref{lem:summary-monitor}). The track flags that keep the column word
valid, the order of the summary updates, and the space accounting are in
\cref{app:complexity}.}
\Description{An algorithm box with an accept move, which tests a short unary word, and a derive move, which streams the columns of a tuple of blocks and replaces the current monitor by the monitor of their summary.}
\label{fig:summary-search}
\end{figure}

\begin{figure}[tbp]
\centering
\begin{tikzpicture}[>=Stealth,
    mon/.style={draw, rounded corners=2pt, align=center, font=\small,
      inner sep=3pt, minimum height=9mm}]
  \node[mon] (w0) {$\mathcal W_0=\mathcal M_0$\\[-1pt]over $\Gamma$};
  \node[mon, right=16mm of w0] (w1) {$\mathcal W_1$\\[-1pt]over $[h_1]$};
  \node[right=9mm of w1] (dots) {$\cdots$};
  \node[mon, right=9mm of dots] (wj) {$\mathcal W_j$\\[-1pt]over $[h_j]$};
  \node[mon, right=12mm of wj] (acc) {accept:\\[-1pt]$z^{m'}$ is a hole};
  \draw[->] (w0) -- node[above,font=\footnotesize] {$\ell_0,\Theta_1$}
    node[below,font=\footnotesize,align=center] {columns\\[-2pt]streamed} (w1);
  \draw[->] (w1) -- node[above,font=\footnotesize] {$\ell_1,\Theta_2$} (dots);
  \draw[->] (dots) -- (wj);
  \draw[->] (wj) -- node[above,font=\footnotesize] {$z,m'$} (acc);
  \draw[->, dashed] ([yshift=-5mm]acc.south) -- node[below,font=\footnotesize]
    {expansion into a hole of $\mathcal M_0$, never stored}
    ([yshift=-5mm]w0.south);
\end{tikzpicture}
\caption{The summary search builds the derived monitors bottom-up. A derive
step chooses a marker $\ell_j$ of the current monitor $\mathcal W_j$, streams
the columns of $h_{j+1}$ distinct blocks that avoid it, keeps only their
summary $\Theta_{j+1}$, and moves to
$\mathcal W_{j+1}=\mathcal W_j[\ell_j,\Theta_{j+1}]$ over the abstract
letters $[h_{j+1}]$. Expanding the abstract letters level by level would
turn an accepted unary word into a hole of $\mathcal M_0$ (dashed arrow).
Compression in \cref{sec:search} instead starts from a given hole and
shortens its hierarchy top-down.}
\Description{A chain of current monitors from the original monitor to a level at which a unary word is accepted, with a dashed arrow back indicating the expansion into a hole that is never stored.}
\label{fig:summary-levels}
\end{figure}

\begin{proposition}[Summary search]\label{prop:summary-search}
Let $\mathcal M_0$ be a lettered monitor with $N\ge1$ states over a finite
alphabet $\Gamma$, given by its tables, and let $d\ge0$. Whether
$\mathcal M_0$ has a nonempty hole of depth at most $d$ is decidable in
nondeterministic space
\[
 O\bigl((N+d+1)^2\,(\DictBound(N,d)+|\Gamma|)^2\bigr),
\]
which is polynomial in $N+d+1$, the dictionary bound $\DictBound(N,d)$ of
\eqref{eq:dictionary-constants}, and $|\Gamma|$. Hence it is decidable in
deterministic space $2^{O((N+d+1)^3)}\cdot(1+|\Gamma|)^4$.
\end{proposition}

The search of \cref{fig:summary-search} has two moves. An accept move tests
a bounded unary hole of the current monitor. A derive move chooses a marker
and streams a bounded tuple of blocks that avoid it. The search maintains
the following invariant. At each level $j\ge1$, each current letter names a
distinct block over the letters of level $j-1$. The stored tables give the
transitions, final targets, and admissions of the corresponding derived
monitor on these letters. A derive move generates the blocks one column at a
time, checks from the equality bits of the summary that they are distinct,
and then retains only their names $1,\ldots,h$ and the tables of
$\mathcal W[\ell,\Theta]$. Higher levels use only this information.

An accepting computation could be expanded backward into a hole of
$\mathcal M_0$: replace each abstract letter by the block it names, level by
level, and close each block with its marker. Agreement
(\cref{lem:agreement-derivative}) and desubstitution (\cref{lem:der-holes})
make each expansion a hole of the monitor one level below, which gives
soundness; the search itself never performs or stores this expansion.
Completeness follows from the dictionary bound (\cref{lem:dictionary-bound}),
column shortening, and unary shortening (\cref{app:search}): the markers and
the shortened blocks of a suitable hole provide the guesses. The largest
stored items are the relations of the current monitor, the derived-admission
bits of one summary, and the column counter, each of
$O((N+d+1)^2(\DictBound(N,d)+|\Gamma|)^2)$ bits. The deterministic bound then
follows from Savitch's theorem~\citep{Savitch70}. The full proof is in
\cref{app:complexity}.

\begin{theorem}[Upper bounds]\label{thm:complexity}\label{cor:complexity}
For every arity $n$, every variance $c$, and each of the three tree domains,
entailment $\phi\models_{\mathcal D,c}x\le y$ of a flat system with $k$
variables is decidable in deterministic space
$2^{2^{O(k)}}\cdot\operatorname{poly}(\insize\phi)$. Hence entailment and
non-entailment are in 2-EXPSPACE and in 3-EXPTIME in each domain, also for
the general input syntax of \cref{prop:flattening}.
\end{theorem}

\begin{proof}[Proof sketch]
By \cref{lem:poly-tests}, decide in polynomial time whether $\phi$ is
satisfiable in $\mathcal D$; if it is not, the entailment holds. Otherwise,
by \cref{prop:transfer}, \eqref{eq:g13}, and \cref{thm:holes}, the entailment
holds iff neither label reader has a hole. The empty word is a hole iff the
initial state is rejected. By \cref{thm:brt,thm:depth}, a label reader has a
hole iff the empty word is a hole or the reader has a nonempty hole of depth
at most $d=8N^2-2N-1$. Its nonempty holes are those of its lettered monitor,
whose tables are computed entry by entry by \cref{lem:poly-tests}.
\Cref{prop:summary-search}, with $|\Gamma|=n$ and
$N+d+1=8N^2-N\le2^{8k+5}$, decides this in space $2^{2^{O(k)}}(1+n)^4$.
Since $k\le\insize\phi$, the whole procedure runs in space
$2^{2^{O(k)}}\operatorname{poly}(\insize\phi)$, hence in time
$2^{2^{2^{O(\insize\phi)}}}$. At positive arity, \cref{prop:flattening}
reduces the general syntax in polynomial time to a flat system with
$O(\insize\phi)$ variables; at arity zero, \cref{prop:arity-zero} in
\cref{app:small-arity} applies.
\Cref{app:complexity} gives the full accounting.
\end{proof}

\paragraph{Comparison with the Length Bound}
Guessing a hole of length at most $H$ and checking it places non-entailment
in 3-NEXPTIME and entailment in co-3-NEXPTIME, hence in 3-EXPSPACE; these are
the bounds that the length bound alone provides. \Cref{thm:complexity}
implies them, since
\[
 \text{2-EXPSPACE}\subseteq\text{3-EXPTIME}\subseteq
 \text{3-NEXPTIME}\cap\text{co-3-NEXPTIME},
\]
and it lowers the space bound by one exponential; by the space hierarchy
theorem \citep[Theorem~6]{StearnsHartmanisLewis65}, 2-EXPSPACE is a proper
subclass of 3-EXPSPACE. The bound $H$ on the hole length is unchanged. The
saving comes from representing one level of the hierarchy by the tables of
its derived monitor instead of storing a hole. For fixed $n\ge2$:
\begin{center}
\small
\begin{tabular}{@{}lll@{}}
\toprule
Procedure & Stored object & Resulting bound\\
\midrule
Enumeration (\cref{fig:algorithm}) & a word of length at most $H$ & time: four exponentials in $k$\\
Guessing a hole & a hole of length at most $H$ & 3-NEXPTIME, hence 3-EXPSPACE\\
Summary search (\cref{prop:summary-search}) & tables of one level & 2-EXPSPACE and 3-EXPTIME\\
\bottomrule
\end{tabular}
\end{center}
The search keeps tables of doubly exponential size in $k$, whereas the only
known bound on the hole length is triply exponential.

% complexity-lower: exponential holes and PSPACE-hardness for arity n >= 2

\subsection{Lower Bounds}
\label{sec:complexity-lower}

We prove two lower bounds for a covariant constructor of arity $n\ge2$:
holes can be forced to be exponentially long in $k$, and entailment is
PSPACE-hard, in each of the three domains.
All systems in this subsection are covariant, so $H_0$ is the identity in
\eqref{eq:g12}. A word $w$ is then left unsafe for $(\phi,x,y)$ iff some
solution $\sigma$ has $\widehat{\sigma(x)}(w)=\top\ne\widehat{\sigma(y)}(w)$,
and right unsafe iff some solution has
$\widehat{\sigma(y)}(w)=\bot\ne\widehat{\sigma(x)}(w)$. Proofs omitted here
are given in \cref{app:lower}.

\paragraph{Overview}
Both constructions simulate a nondeterministic automaton over $\{0,1\}$ by
constructor literals along binary paths. For long holes, lower literals run
an automaton on $y$: a transition $q\xrightarrow{a}q'$ places the $a$-child
of $Y_q$ above $Y_{q'}$, so $Y_q$ has no bottom prefix along any word that
the automaton accepts from $q$. All states of the automaton are final, so
its language is prefix closed, and the remaining words form a union of sets
$u\{0,1\}^*$. We use an automaton that accepts every word except the
extensions of one exponentially long word, the record of a binary counter.
Separate upper literals keep $x$ top free, with constructor positions in
$\{0,1\}^*$, so every unsafe word is right unsafe and extends the counter
record. Truncating the trees of the automaton states, with top leaves at a
depth beyond a chosen word, gives finite solutions, so these words are
unsafe in all three domains. For hardness the roles of the two kinds of
literals are reversed: upper literals run a given automaton on $x$, and lower
literals keep a second variable $r$ at least a constructor along every
proper prefix of a code word.

\paragraph{Constructor Positions}
For a tree $t$, let $\Pos(t)=\{w\in\operatorname{dom}(t):t(w)=f\}$ be its
set of \emph{constructor positions}; it is prefix closed. A tree is
\emph{top free} if no node is labeled $\top$. Every prefix-closed set
$P\subseteq\Dir^*$ is $\Pos(T_P)$ for exactly one top-free tree $T_P$: its
domain is $\{\eps\}\cup P\cup P\Dir$, its label is $f$ on $P$, and $\bot$
elsewhere. For example, $T_\varnothing=\Bot$, and $T_{\{0,1\}^*}$ is the
full constructor tree $F$ when $n=2$. For a prefix-closed set $P$ and
$m\ge1$, the \emph{truncation} $T_P^m$ is obtained from $T_P$ by replacing the
subtree at every word of $P$ of length $m$ by $\Top$; it is a finite tree.
For variables $a_0,a_1,b$ we write $f(a_0,a_1,\bar b)$ for
$f(a_0,a_1,b,\ldots,b)$, with $n-2$ copies of $b$; for $n=2$ it stands for
$f(a_0,a_1)$.

\begin{lemma}[Top-free trees]\label{lem:cx-topfree}
Let $n\ge2$.
\begin{enumerate}[leftmargin=*]
\item For top-free trees $t,u$ and every word $w$, $\widehat t(w)=f$ if
  $w\in\Pos(t)$ and $\widehat t(w)=\bot$ otherwise; and $t\le u$ iff
  $\Pos(t)\subseteq\Pos(u)$.
\item $T_P$ is finite iff $P$ is finite.
\item Let $\sigma$ satisfy $b=\Bot$ and $g\le f(g,g,\bar b)$. Then
  $\sigma(g)$ is top free and $\Pos(\sigma(g))\subseteq\{0,1\}^*$, and so is
  every tree $t\le f(\sigma(g),\sigma(g),\Bot,\ldots,\Bot)$.
\item Let $P,P'$ be prefix closed, $\eps\in P$, $a\in\Dir$, and
  $aP'\subseteq P$. Then $f(t_0,\ldots,t_{n-1})\le T_P^m$ for every $m\ge1$,
  where $t_a=T_{P'}^m$ and $t_i=\Bot$ for $i\ne a$.
\end{enumerate}
\end{lemma}

Item~(3) holds because a top leaf cannot lie below a constructor: by
\eqref{eq:g11}, a top leaf of $\sigma(g)$ at a word $iw$ would force
$i\in\{0,1\}$ and a top leaf at the shorter word $w$, and $\sigma(g)$ has no
top leaf at $\eps$. In item~(4), the truncation at the same depth $m$ on both
sides is what makes the inequality hold: a top leaf of $T_{P'}^m$ at a word
$v$ of length $m$ lies below the top leaf of $T_P^m$ at the prefix of $av$ of
length $m$.

\paragraph{Long Holes}
As noted in \cref{sec:complexity-upper}, a failed entailment on a
satisfiable system is certified by a hole of one of the two label readers,
and by \cref{thm:holes} these holes are exactly the unsafe words. Such
certificates can be forced to be exponentially long. The construction uses
an automaton with $O(\beta)$ states for the binary words that do not begin
with one word of length $2(\beta+1)2^\beta$, the record of a $\beta$-bit
binary counter.

Fix $\beta\ge1$. Over the \emph{token alphabet} $\mathbb T=\{0,1,\#,\$\}$,
let $\cxbin j$ be the binary expansion of $j$ with $\beta$ bits, least
significant bit first, and put
\[
 \mathsf H_\beta=\cxbin0\#\cxbin1\#\cdots\#\cxbin{2^\beta-1}\$,
 \qquad \len{\mathsf H_\beta}=(\beta+1)2^\beta.
\]
The \emph{token encoding} maps the letters of $\mathbb T$ to $0\mapsto00$,
$1\mapsto01$, $\#\mapsto10$, and $\$\mapsto11$. Let
$\mathsf W_\beta\in\{0,1\}^*$ be the token encoding of $\mathsf H_\beta$, of
length $\lambda_\beta=2(\beta+1)2^\beta$. For $\beta=1$,
$\mathsf H_1=0\#1\$$ and $\mathsf W_1=00100111$. In a nondeterministic
automaton all of whose states are final, a word is accepted iff it has a
run, that is, a path labeled by the word from an initial state.

\begin{lemma}[Counter automaton]\label{lem:cx-counter}
For every $\beta\ge1$ there is a nondeterministic automaton $\mathcal B_\beta$
over $\{0,1\}$ with $9\beta+27$ states, all of them final, whose language is
$\{0,1\}^*\setminus\mathsf W_\beta\{0,1\}^*$. It is computable in time
polynomial in $\beta$.
\end{lemma}

The automaton reads the encoded tokens. In parallel, its branches check the
format of the blocks, that the first block is zero, that each block is the
successor of the previous one, and that $\$$ follows exactly the all-ones
block. A branch that detects an error enters a state that loops on every
letter, and on $\mathsf W_\beta$ itself every branch dies (\cref{app:lower}).
For instance, for $\beta=2$ the record is $\mathsf H_2=00\#10\#01\#11\$$. In a
token word that begins with $00\#10\#11$, the first bit of the third block is
wrong: the successor of $10$, that is, of $1$, is $01$. The branch spawned at
the first bit of the second block predicts the bit $0$ from that bit and the
carry, counts $\beta+1=3$ tokens, and reads $1$ instead. It detects the
error while storing only the predicted bit and a counter, not the history.

Let $Q$ be the set of states of $\mathcal B_\beta$, $I\subseteq Q$ its
initial states, and $q\xrightarrow{a}q'$ its transitions. Add a state
$q_{\mathrm{in}}$ with $q_{\mathrm{in}}\xrightarrow{a}q'$ iff
$q\xrightarrow{a}q'$ for some $q\in I$, so that $q_{\mathrm{in}}$ accepts the
language of $\mathcal B_\beta$. For $n\ge2$, let $\Phi_\beta^{(n)}$ have the
variables $x$, $g$, $b$, and $Y_q$ for $q\in Q\cup\{q_{\mathrm{in}}\}$, with
$y=Y_{q_{\mathrm{in}}}$, and the flat literals
\[
\begin{aligned}
 &b=\Bot,\quad g\le f(g,g,\bar b),\quad x\le f(g,g,\bar b)
   &&\text{(confinement to binary paths)},\\
 &f(b,b,\bar b)\le Y_q &&\text{(nonbottom roots)},\\
 &f(Y_{q'},b,\bar b)\le Y_q\ \ (q\xrightarrow0q'),\quad
  f(b,Y_{q'},\bar b)\le Y_q\ \ (q\xrightarrow1q')
   &&\text{(propagation along transitions)},
\end{aligned}
\]
for all $q\in Q\cup\{q_{\mathrm{in}}\}$. Thus $y$ receives the lower literals
of the initial states directly. Write $\Phi_\beta=\Phi_\beta^{(2)}$. For every
$n$, $\Phi_\beta^{(n)}$ has $k=9\beta+31$ variables. The confining literals
keep $x$ top free, with constructor positions in $\{0,1\}^*$. The other two
groups give each $Y_q$ a constructor or a top prefix along every word that
$\mathcal B_\beta$ accepts from $q$, so $y$ has no bottom prefix along any
binary word without the prefix $\mathsf W_\beta$. The trees of $y$ and of the $Y_q$
may have top leaves, which permits finite solutions.

\begin{proposition}[Exponentially long holes]\label{prop:cx-length}
Let $n\ge2$, let $f$ be covariant of arity $n$, let $\beta\ge1$, and let
$\Phi=\Phi_\beta^{(n)}$, with the query $x\le y$.
\begin{enumerate}[leftmargin=*]
\item $\Phi$ has a finite solution, and $\Phi\models_{\mathcal D}x\le y$ fails
  in each of the three domains $\mathcal D$.
\item No word is left unsafe, and the right-unsafe words are exactly the
  words of $\mathsf W_\beta\{0,1\}^*$.
\item The left label reader
  $(\mathcal M_{\mathsf l},\Adm_{\mathsf l},\Rej_{\mathsf l})$ has no hole,
  and the holes of the right label reader
  $(\mathcal M_{\mathsf r},\Adm_{\mathsf r},\Rej_{\mathsf r})$ are exactly the
  words of $\mathsf W_\beta\{0,1\}^*$. In particular, every hole that the
  search of \cref{fig:algorithm} can find has length at least
  $\lambda_\beta=2(\beta+1)2^\beta$, and $\mathsf W_\beta$ is a hole of this
  length.
\end{enumerate}
Hence, for every fixed arity $n\ge2$ and in each domain, shortest holes of
systems with $k$ variables can have length $2^{\Omega(k)}$: here
$\lambda_\beta\ge2^{(k-22)/9}$.
\end{proposition}

\begin{proof}[Proof idea]
For $n=2$, let $\sigma$ be a solution. By \cref{lem:cx-topfree}(3), $\sigma(x)$
is top free with $\Pos(\sigma(x))\subseteq\{0,1\}^*$, so no word is left
unsafe, and a right-unsafe word $w$ has $w\in\Pos(\sigma(x))$ and
$\widehat{\sigma(y)}(w)=\bot$ for some solution $\sigma$. Induction along runs of
$\mathcal B_\beta$ shows that $\widehat{\sigma(y)}$ is not $\bot$ on the language
of $\mathcal B_\beta$, so by \cref{lem:cx-counter} every right-unsafe word
begins with $\mathsf W_\beta$. Conversely, let $w\in\mathsf W_\beta\{0,1\}^*$
and $m=\len w+1$. The truncations $Y_q=T_{L_q}^m$, with $L_q$ the set of words
accepted from $q$, together with $g=T_{\{0,1\}^{<m}}$, $b=\Bot$, and $x=T_P$ for
the set $P$ of prefixes of $w$, form a finite solution with $w\in\Pos(x)$ and
$\widehat y(w)=\bot$ (\cref{lem:cx-topfree}(4)). This gives~(1) and~(2),
and~(3) follows by \cref{thm:holes}. The case $n\ge3$ reduces to $n=2$ by an
embedding of binary trees, described at the end of this subsection. The
details are in \cref{app:lower}.
\end{proof}

Since $N=2^{4k+1}$, the length $\lambda_\beta$ is only polynomial in $N$,
whereas the length bound $H$ of \cref{sec:complexity-upper} is triply
exponential in $k$.

\paragraph{Hardness}
\Citet[Theorems~1 and~2]{HengleinRehof98} proved PSPACE-hardness for a binary
covariant constructor over arbitrary and over finite trees
(\cref{sec:related}). We give a reduction from universality of
nondeterministic automata over $\{0,1\}$ that works in all three domains and
for every fixed arity $n\ge2$. An automaton $\mathcal A=(Q,q_0,F,\delta)$ has
one initial state $q_0$, final states $F$, and transitions
$\delta\subseteq Q\times\{0,1\}\times Q$; it is universal if it accepts every
word. Universality of such automata is PSPACE-hard.\footnote{%
\citet[Theorem~3]{KrotzschMasopustThomazo16} prove this for automata over
$\{0,1\}$, even for partially ordered ones. Their automata may have several
initial states; a fresh initial state that has the outgoing transitions of
all of them, and is final iff one of them is, accepts the same language. For
regular expressions over $\{0,1\}$ with union, concatenation, and star, the
classical result is due to \citet[Section~3]{StockmeyerMeyer73}, who state
that inequivalence is complete for nondeterministic linear space.}

The reduction runs an automaton on $x$ by upper literals, and an accepting
run forces the subtree of $x$ at its end to be $\Bot$; then $x$ has no node
below that word. A word at which $x$ is to have a constructor while a second
variable $r$ is $\bot$ must therefore have no accepted prefix. We encode
words by a prefix-free code, so that the code of a rejected word has no
accepted code as a prefix. The variables $r$ and $r'$ track the position
within an encoded letter (\cref{fig:cx-code}). Their lower literals keep $r$
at least a constructor at every proper prefix of a code, while at a
completed code $r$ may be $\bot$, where $x$ can still have a constructor if
the code is rejected.

For $w\in\{0,1\}^*$ let $\cxcode(w)=e(w)1$, where $e$ replaces each letter
$0$ by $00$ and each $1$ by $01$; we call $\cxcode$ the \emph{prefix-free
code}. On single bits, $e$ agrees with the token encoding, but the two
encodings serve different constructions. The code is prefix
free: if $\cxcode(w')\pre\cxcode(w)$ and $\len{w'}<\len w$, then position
$2\len{w'}$ carries the final $1$ of $\cxcode(w')$ but the letter $0$ in
$\cxcode(w)$. From $\mathcal A$ we obtain in linear time an automaton
$\mathcal A'$, again with the initial state $q_0$, that accepts
$\cxcode(L(\mathcal A))$: replace each transition $(q,a,q')$ by
$q\xrightarrow0\mu\xrightarrow{a}q'$ through a new middle state $\mu$, and
add a transition labeled $1$ from every final state to one new state
$q_{\mathrm f}$, the only final state of $\mathcal A'$.

For such an automaton $\mathcal A'$ with states $Q'$, let $\chi(\mathcal A')$
have the variables $z$, $b$, $r$, $r'$, and $X_q$ for $q\in Q'$, write $x$ for
$X_{q_0}$, and take the flat literals
\[
\begin{gathered}
 X_q\le f(X_{q'},z,b,\ldots,b)\ \ (q\xrightarrow0q'),\qquad
 X_q\le f(z,X_{q'},b,\ldots,b)\ \ (q\xrightarrow1q'),\qquad
 X_{q_{\mathrm f}}=\Bot,\\
 b=\Bot,\qquad f(r',b,\ldots,b)\le r,\qquad f(r,r,b,\ldots,b)\le r'.
\end{gathered}
\]
Here $b$ fills the arguments $2,\ldots,n-1$ of the first two literals and
the arguments $1,\ldots,n-1$ of the first lower literal; for $n=2$ the first
two literals have no $b$. No literal bounds $z$, so $z$ may be $\top$, and
then, among the two binary children, an upper literal constrains only the child
that the run follows; the added children are bounded by $b=\Bot$. Along
every run of $\mathcal A'$ from $q_0$ on a word $u$, the upper literals keep the subtree
of $x$ at $u$, if it exists, below the variable $X_q$ of the reached state
$q$; an accepting run therefore forces the subtree of $x$ at its end, if it
exists, to be $\Bot$. The lower literals give $r$ a constructor or $\top$ at
every proper prefix of a code, and they allow $r$ to be $\bot$ at a code. If
$\mathcal A$ is universal, every proper prefix of a code lies on an accepting
run, and $x\le r$ follows. If $\mathcal A$ rejects a word $w$, then a solution
can give $x$ a constructor at $\cxcode(w)$ and $r$ the label $\bot$ there.

\begin{figure}[tbp]
\centering
\begin{tikzpicture}[>=Stealth, auto,
    every state/.style={minimum size=9mm, inner sep=1pt}]
  \node[state, initial, initial text={}] (r) at (0,0) {$r$};
  \node[state] (rp) at (3.2,0) {$r'$};
  \node[draw, rounded corners=2pt, font=\small, align=center] (bot) at (0,-1.9)
    {bottom permitted\\[-1pt]at a completed code};
  \draw[->] (r) edge[bend left=25] node {$0$} (rp);
  \draw[->] (rp) edge[bend left=25] node {$0,1$} (r);
  \draw[->, dashed] (r) -- node[left] {$1$} (bot);
\end{tikzpicture}
\caption{The lower literals $f(r',b,\ldots,b)\le r$ and
$f(r,r,b,\ldots,b)\le r'$ of $\chi(\mathcal A')$ as a diagram: an edge
labeled $a$ from $v$ to $v'$ means that the $a$-child of $v$ lies above $v'$.
Reading the encoding $e(w)$ of a word follows the solid edges, so the tree
of $r$ has a constructor or $\top$ at every proper prefix of a code. The
final $1$ of $\cxcode(w)$ leaves the diagram (dashed): at a completed code the
tree of $r$ may have the label $\Bot$, although $r$ itself is never $\Bot$.}
\Description{Two states r and r prime with a 0-edge from r to r prime and a 0,1-edge back, and a dashed 1-edge from r to a box saying that the tree of r may be bottom at a completed code.}
\label{fig:cx-code}
\end{figure}

\begin{theorem}[PSPACE-hardness]\label{thm:cx-hard}
Let $n\ge2$ be fixed and let $f$ be covariant. For every automaton
$\mathcal A$ over $\{0,1\}$ and each domain $\mathcal D$,
\[
 \chi(\mathcal A')\models_{\mathcal D}x\le r
 \quad\Longleftrightarrow\quad
 \mathcal A\text{ is universal}.
\]
Consequently, in each of the three domains, entailment of flat systems over
$\{\Bot,f,\Top\}$ is PSPACE-hard. So is the problem with the variance as part
of the input, already on covariant instances.
\end{theorem}

\begin{proof}[Proof idea]
Let $P$ be the set of proper prefixes of codes, the words $e(w)$ and
$e(w)0$. If $\mathcal A$ is universal, every word of $P$ lies on an accepting
run of $\mathcal A'$, so in every solution the subtree of $x$ there, if
present, lies below a constructor tree whose arguments $2,\ldots,n-1$ are
$\Bot$; every code is accepted, so $x$ is $\Bot$ at codes. Along $P$, the
lower literals give $r$ only the labels $f$ and $\top$. Every common position
of $x$ and $r$ therefore compares a label at most $f$ with one at least $f$,
or has $x=\Bot$. If $\mathcal A$ rejects $w$, no prefix of $\cxcode(w)$ is
accepted by $\mathcal A'$, and a path tree along $\cxcode(w)$ together with
finite approximants of $r$ forms a finite countermodel, which lies in every
domain. The system is satisfiable, has $\len{Q'}+4$ variables, and is
computed in polynomial time. The details are in \cref{app:lower}.
\end{proof}

For $n=2$ and arbitrary or finite trees, \cref{thm:cx-hard} reproves, by a
different reduction, the hardness theorems of
\citet[Theorems~1 and~2]{HengleinRehof98}, which are stated for a binary
covariant constructor over these two domains. It gives a uniform proof for
all three domains and every fixed arity $n\ge2$. Their theorems do not state
the regular case, and we do not claim that it is new. We claim no hardness
for a fixed variance with a contravariant argument, nor for arity one
(\cref{app:small-arity}).

\paragraph{Larger Arities}
The hardness reduction works directly at every arity $n\ge2$. The systems
$\Phi_\beta^{(n)}$ for long holes are obtained from $\Phi_\beta$ by filling
the new arguments of every constructor with the variable $b=\Bot$. Such
\emph{bare padding} does not preserve lower bounds in general. The padded
system also constrains trees that are not images of binary trees, and these
may refute the entailment at short words that use the new letters.

\begin{example}[Bare padding fails]\label{ex:cx-padding}
For a binary covariant $f_2$, the system
\[
 \chi=\{x\le f_2(y,b),\ f_2(z,b)\le y,\ f_2(z,t)\le z,\ b=\Bot,\ t=\Top\}
\]
entails $x\le y$. Indeed, every top leaf of $x$ lies at a word $0v$ such that
$y$ has a top prefix along $v$. Then $y$ also has one along $0v$: either
$y=\Top$, or the child $0$ of $y$ lies above $z$, which has a top prefix along
every word that begins with $0^j1$; and for $v\in0^*$, the top prefix along
$v$ lies along $0v$. Moreover, $y$ has bottom prefixes only along words that
begin with $1$, where $x$ has one as well. Padding with a variable $o=\Bot$
in the third argument, for $n=3$, loses the entailment: the finite assignment
$x=f(f(\Bot,\Bot,\Top),\Bot,\Bot)$, $y=f(z,\Bot,\Top)$, $z=f(\Top,\Top,\Bot)$,
$b=o=\Bot$, and $t=\Top$ is a countermodel, since $x$ has a top leaf at $02$,
where $y$ has the bottom leaf of $z$.
\end{example}

Padding with $\Bot$ does preserve solutions and unsafe words on binary
paths: lifting a binary tree by adding $\Bot$ children, and projecting a tree
over $\Dir$ to its binary paths, map solutions to solutions in each domain
and keep leaf extensions on $\{0,1\}^*$ (\cref{prop:cx-embedding} in
\cref{app:lower}). For $\Phi_\beta^{(n)}$, the confining literals
$g\le f(g,g,\bar b)$ and $x\le f(g,g,\bar b)$ also exclude unsafe words that
use an added child position, because they keep $x$ top free with constructor
positions in $\{0,1\}^*$. Thus $\Phi_\beta^{(n)}$ and $\Phi_\beta$ have
exactly the same unsafe words. The system of \cref{ex:cx-padding} has no
confining literals, and the word $02$ of its countermodel is left unsafe and
uses the added letter $2$. The embedding is therefore not a general
reduction between arities, such as the equivalence of all arities at least
two asserted by \citet[Corollary~1]{NiehrenPriesnitz03} (\cref{sec:errata}).
It transfers unsafe words only on binary paths, which suffices when
confining literals exclude all other words.

\paragraph{Combined Bounds}
Together with \cref{thm:complexity}, entailment for a covariant constructor
of each fixed arity $n\ge2$ is PSPACE-hard and in 2-EXPSPACE in each domain.
\Cref{sec:future} discusses the open questions about its exact complexity.

